\chapter{Preservation Mode, Reconstitution Mode, and the Fidelity Continuum}
\label{ch:fidelity-continuum}

\begin{quotation}
\noindent\textit{Synopsis.} This chapter revisits the traditional
distinction between direct access and reconstruction and replaces it with a
single fidelity continuum. Every consumer interacts with some rendering,
and every rendering preserves only part of the original distinction
structure, so preservation becomes a limiting case of reconstitution rather
than a separate ontology. Fidelity is introduced as a measure of how
closely a rendering preserves an original object's admissible continuation
space, and the chapter argues that reconstitution does not escape the
Composition Failure Theorem: a reconstructed trace is simply a new relay
stage, subject to the same constraints as any other.
\end{quotation}

\section{The Naive Binary}

An initial formulation of access modes distinguishes preservation mode --
the consumer has direct causal access to the artifact itself -- from
reconstitution mode -- the consumer has access only to traces from which
structure must be regenerated. This chapter shows that the binary does
not survive scrutiny and replaces it with a single continuous variable.

\begin{remark}
\label{rem:binary-fails}
Nearly every case initially classified as preservation is in fact
mediated by some rendering: a photocopy, a scan, a decoded audio file. A
photograph of a painting is not the painting; a scan of a manuscript is
not the manuscript; a decoded audio stream is not the acoustic event it
encodes. Such cases are already points on the loss spectrum of
\cref{def:rendering}, not a qualitatively separate category.
\end{remark}

\section{Fidelity, Defined}

\begin{definition}[Fidelity]
\label{def:fidelity}
For a rendering $r$ of an object $e$ and a task $\tau$,
\[
\fidelity{r}{\tau} = \frac{\lvert \admis{\tau}{r} \cap \admis{\tau}{e}
\rvert}{\lvert \admis{\tau}{e} \rvert},
\]
a measure of how closely the admissible continuation set available from
$r$ approximates the admissible continuation set available from the
original object $e$ itself, for the task $\tau$ in question. Per
\cref{rem:finite-distinctions}, this is stated only for the finite or
countable case.
\end{definition}

\begin{remark}[Relation to Transport Stability]
\label{rem:fidelity-transport-pointer}
Fidelity is not the only available measure of how well structure survives a
transformation. Distinguishability Geometry's Transport Stability theorem
bounds a different but related quantity, and Chapter~\ref{ch:four-operations}
(\cref{rem:fidelity-transport-open}) states plainly that this monograph has
not shown one reduces to the other; the relationship is recorded as open
rather than assumed.
\end{remark}

\begin{remark}[The Intersection Simplifies]
\label{rem:fidelity-simplifies}
By \cref{prop:admissibility-upper-bound}, $\admis{\tau}{r} \subseteq
\admis{\tau}{e}$ always, so the intersection in \cref{def:fidelity} is
simply $\admis{\tau}{r}$ itself:
$\fidelity{r}{\tau} = \lvert\admis{\tau}{r}\rvert / \lvert\admis{\tau}{e}\rvert$.
The intersection form is retained in the definition because it makes
explicit, without relying on a separately proved proposition, exactly
what is being compared; the simplification is used freely from here on.
\end{remark}

\begin{remark}[A Standing Assumption: Ground-Truth Sufficiency]
\label{rem:ground-truth-sufficiency}
Throughout this chapter it is assumed that $\admis{\tau}{e} \subseteq
\mathcal{C}_\tau(e)$: the true object $e$ is itself always sufficient for
any task evaluated against it. This is close to definitional given how
$\mathcal{C}_\tau(e)$ was fixed relative to $e$ in
\cref{def:correctness-set}, but it is worth stating openly, since
\cref{cor:fidelity-implies-sufficiency} below depends on it.
\end{remark}

\begin{proposition}[Preservation as a Limit Case]
\label{prop:preservation-limit}
``Preservation mode'' is the limit of reconstitution mode as
$\fidelity{r}{\tau} \to 1$ for the tasks $\tau$ in question; it is not a
distinct ontological category but a point on the continuum defined by
\cref{def:fidelity}.
\end{proposition}

\begin{proofsketch}
As $\fidelity{r}{\tau} \to 1$, $\admis{\tau}{r}$ converges to
$\admis{\tau}{e}$: the continuations available from the rendering become
indistinguishable, for the purposes of $\tau$, from those available from
the original. Reconstitution from a sufficiently high-fidelity trace is
then operationally indistinguishable from direct access, which is what
``preservation'' was intended to capture.
\end{proofsketch}

\begin{corollary}[Fidelity One Implies Sufficiency]
\label{cor:fidelity-implies-sufficiency}
If $\fidelity{r}{\tau} = 1$, then $r$ is sufficient for $\tau$
(\cref{def:sufficiency}).
\end{corollary}

\begin{proofsketch}
$\fidelity{r}{\tau}=1$ gives $\admis{\tau}{r} = \admis{\tau}{e}$ by
\cref{rem:fidelity-simplifies}, and $\admis{\tau}{e} \subseteq
\mathcal{C}_\tau(e)$ by \cref{rem:ground-truth-sufficiency}, so
$\admis{\tau}{r} \subseteq \mathcal{C}_\tau(e)$, which is exactly
\cref{def:sufficiency}.
\end{proofsketch}

This corollary sharpens \cref{prop:preservation-limit}'s informal
``operationally indistinguishable'' into a precise claim: the fidelity
continuum's upper limit is not merely subjectively hard to tell apart
from direct access, it is provably sufficient in the sense
Chapter~\ref{ch:admissibility} already uses throughout the rest of the
monograph.

\begin{proposition}[Benign Loss Preserves Perfect Fidelity]
\label{prop:benign-preserves-fidelity}
If every distinction in $\discard{r}$ is benign relative to $\tau$
(\cref{def:benign-loss}), then $\fidelity{r}{\tau} = 1$.
\end{proposition}

\begin{proofsketch}
By the same argument as \cref{prop:benign-preserves-sufficiency}: no
distinction discarded by $r$ is required to determine any continuation's
membership in $\admis{\tau}{e}$, so no continuation reachable from $e$
becomes unreachable from $r$ for the purposes of $\tau$, giving
$\admis{\tau}{r} = \admis{\tau}{e}$ and hence
$\fidelity{r}{\tau} = 1$ by \cref{rem:fidelity-simplifies}.
\end{proofsketch}

\cref{prop:benign-preserves-fidelity} is the fidelity-continuum
counterpart of \cref{prop:benign-preserves-sufficiency}, and the two
propositions together are why this chapter's continuous measure and
Chapter~\ref{ch:admissibility}'s binary sufficiency criterion cohere
rather than compete: perfect fidelity for a task and sufficiency for that
task are the same claim, viewed through cardinality in one case and
through set containment in the other.

\section{The Hard Test Case: A Chess Board in Front of You}

A chess board physically present to a player is the strongest available
candidate for genuine, unmediated preservation. Yet perceiving the board
is itself a rendering process -- a visual system producing a percept from
ambient light -- and is therefore, by \cref{rem:binary-fails}, already a
point on the fidelity continuum rather than a case falling outside it.
Whether \emph{any} case achieves fidelity exactly equal to $1$, or
whether the entire category of strictly unmediated access is an
idealization available only to formal or machine systems (a verifier
checking a proof directly against explicit axioms, per
Chapter~\ref{ch:proofs}) and never to embodied perception, is left open
here.

\begin{openquestion}
\label{oq:unmediated-access}
Does any instance of ``direct access'' to an artifact involve strictly
zero mediation, or is unmediated access an idealization that applies
cleanly only to formal/machine verification and never to embodied
perception? Resolving this requires engaging directly with the question
raised in full in Chapter~\ref{ch:mem8}.
\end{openquestion}

\section{Recursive Insufficiency}

\begin{proposition}[Recursive Insufficiency]
\label{prop:recursive-insufficiency}
If $A_i$ is a trace consumed in reconstitution mode, the act of
reconstitution -- schema plus trace producing a usable object -- is
itself a transition $B_i$ in a budget relay in the sense of
\cref{def:substrate-chain}, and is therefore subject to
\cref{thm:composition-failure}.
\end{proposition}

\begin{proofsketch}
Reconstitution takes a trace $A_i$ and a schema and produces a usable
object $A_{i+1}$ for continuation. This is not merely analogous to a
relay transition; it can be given the same triple structure directly.
The schema's capacity to correctly fill a given gap plays the role of
$C_i$ (some gaps are cheaply and reliably filled by strong priors, others
are not); competition among candidate fillings for the same gap, or
among several gaps drawing on the same limited attention or working
memory, plays the role of $I_i$; and whatever bounded capacity the
reconstitution process has available -- time, attention, cognitive
resources -- plays the role of $\beta_i$. Under this identification,
$A_i \to A_{i+1}$ satisfies \cref{def:substrate-chain} exactly, not by
loose analogy, and \cref{lem:exhaustion}'s mechanism applies to
reconstitution precisely as it applies to any other transition: a
required feature of the reconstructed object that costs more than the
reconstitution process's available capacity is not recovered, regardless
of how the schema is implemented.
\end{proofsketch}

This corrects a claim made informally in the research program that
generated this monograph, in which reconstitution appeared briefly to
``escape'' the insufficiency problem by changing register from
information recovery to regeneration. \cref{prop:recursive-insufficiency}
shows this is not the case, and shows it by identification rather than
metaphor: reconstitution relocates the insufficiency problem to the trace
and the schema, and \cref{thm:composition-failure} applies to that
relocation exactly as it applies to any other stage of a relay, because
it \emph{is} another stage of a relay.
